% !TEX program = xelatex
%% ===========================================================================
%%  main.tex — Laporan Magang
%%  Program Studi Sistem Informasi — STMIK Kuwera
%%
%%  Kompilasi : make            (XeLaTeX via latexmk, sekaligus BibTeX)
%%  Identitas : config/metadata.tex
%%  Gaya      : laporan.sty  (Pedoman Bab 4.1 dan 4.3)
%% ===========================================================================
\documentclass[12pt,a4paper,oneside]{report}

\usepackage{laporan}
%% ===========================================================================
%%  config/metadata.tex — identitas laporan
%%
%%  Ubah HANYA di file ini. Sampul, halaman judul, lembar pengesahan, dan
%%  kata pengantar mengambil data dari sini.
%% ===========================================================================

%% -------------------------------------------------- judul (10-14 kata)
%% Bentuk piramida terbalik: pecah baris manual dengan \\ bila perlu.
%% Judul versi satu baris (tanpa pemisah baris) untuk metadata PDF laporan.
\newcommand*{\JudulLaporanPDF}{Perancangan dan Implementasi Sistem Pencarian Data Riset Berbantuan AI pada Broadband Multimedia Wireless Laboratory, NTUST}

%% Judul sampul (15 kata) dengan pemisah baris membentuk piramida terbalik.
\newcommand*{\JudulLaporan}{Perancangan dan Implementasi Sistem Pencarian Data\\Riset Berbantuan AI pada Broadband Multimedia\\Wireless Laboratory, NTUST}

%% ------------------------------------------------------------- penulis
\newcommand*{\NamaPenulis}{Samuel}   % tanpa gelar dan tanpa singkatan
\newcommand*{\NIM}{2231010010}
\newcommand*{\NIMLine}{NIM \NIM}     % baris NIM di sampul
\newcommand*{\ProgramStudi}{Sistem Informasi}
\newcommand*{\NamaKampus}{STMIK Kuwera}
\newcommand*{\LogoKampus}{assets/stmik_kuwera_logo.png}

%% --------------------------------------------------------------- sampul
\newcommand*{\Kota}{Jakarta}
\newcommand*{\Tahun}{2026}   % tahun laporan dipertahankan dan dinyatakan lulus

%% ---------------------------------------------------- lembar pengesahan
%% Nama setiap penanda tangan dipisah dari gelarnya. Alasannya: baris nama
%% pada lembar pengesahan bergaris bawah dan tidak boleh pindah baris.
\newcommand*{\NamaPembimbing}{Dr. Lulu Setiawati}
\newcommand*{\GelarPembimbing}{S.E., M.Bus., D.Th., CFP, CSA}
\newcommand*{\NIDNPembimbing}{NIDN 0323087303}
\newcommand*{\NamaKetuaProgramStudi}{D. Dedy Prasetya Kristiadi}
\newcommand*{\GelarKetuaProgramStudi}{S.Kom., M.Kom.}
\newcommand*{\NIDNKetuaProgramStudi}{NIDN 0311127402}
\newcommand*{\TanggalPengesahan}{................ 2026}
\newcommand*{\InstansiMagang}{Broadband Multimedia Wireless Laboratory, NTUST}
\newcommand*{\PeriodeMagang}{12 Januari -- 28 September 2026}

%% ------------------------------------------------------- kata pengantar
\newcommand*{\TanggalKataPengantar}{................ 2026}


\hypersetup{%
  pdftitle   = {\JudulLaporanPDF},%
  pdfauthor  = {\NamaPenulis},%
  pdfsubject = {Laporan Magang \ProgramStudi{} \NamaKampus}%
}

\begin{document}

%% --------------------------------------------------- BAGIAN AWAL (Romawi)
%% Sampul dan Halaman Judul: tanpa nomor halaman dan tidak dihitung.
%% Keduanya memang memakai isi yang sama (Pedoman Bab 4.3: "Halaman Judul
%% (isi sama dengan sampul)"): Sampul dicetak di kertas hijau daun sebagai
%% bagian luar, Halaman Judul di kertas biasa sebagai halaman pertama di
%% dalam jilidan. Bila prodi Anda hanya meminta satu halaman sampul,
%% hapus/komentari baris %% ===========================================================================
%%  frontmatter/halaman-judul.tex — Halaman Judul
%%  Isinya sama dengan sampul (Pedoman Bab 4.3), tetapi dicetak di kertas
%%  biasa. Halaman ini juga tidak diberi nomor dan tidak dihitung.
%% ===========================================================================
\BuatSampul

 di bawah.
%% (Angka Romawi besar dipakai di sini hanya agar penanda halaman PDF tidak
%%  bentrok dengan angka Romawi kecil milik bagian awal.)
\pagenumbering{Roman}
%% ===========================================================================
%%  frontmatter/sampul.tex — Sampul
%%  Semua rata tengah, spasi 1. Dicetak pada kertas Omega hijau daun, tinta
%%  hitam. Isi diambil dari config/metadata.tex; tata letak dari laporan.sty.
%% ===========================================================================
\BuatSampul


%% ===========================================================================
%%  frontmatter/halaman-judul.tex — Halaman Judul
%%  Isinya sama dengan sampul (Pedoman Bab 4.3), tetapi dicetak di kertas
%%  biasa. Halaman ini juga tidak diberi nomor dan tidak dihitung.
%% ===========================================================================
\BuatSampul



%% Angka Romawi dihitung mulai Lembar Pengesahan: i, ii, iii, ...
\pagenumbering{roman}
%% ===========================================================================
%%  frontmatter/lembar-pengesahan.tex — Lembar Pengesahan
%%  Halaman pertama yang diberi angka Romawi: i.
%%  Penanda tangan: dosen pembimbing dan Ketua Program Studi.
%% ===========================================================================
\babfrontmatter{Lembar Pengesahan}

\begin{singlespace}
Laporan Magang dengan judul \emph{\JudulLaporanPDF} ini disusun oleh:

\vspace{1\baselineskip}
\begin{tabular}{@{}l@{\hspace{0.5em}}c@{\hspace{0.5em}}l@{}}
Nama          & : & \NamaPenulis \\
NIM           & : & \NIM \\
Program Studi & : & \ProgramStudi \\
\end{tabular}

\vspace{1\baselineskip}
% TODO: sesuaikan kalimat pengantar pengesahan dengan format prodi.
Telah diperiksa dan disetujui untuk dipertanggungjawabkan sebagai laporan
pelaksanaan magang pada \InstansiMagang{} periode \PeriodeMagang.

\vspace{2\baselineskip}
\begin{center}\Kota, \TanggalPengesahan\end{center}

\vspace{2\baselineskip}
\begin{center}
\begin{minipage}[t]{0.6\textwidth}\centering
  Pembimbing,\\[4\baselineskip]
  \uline{\NamaPembimbing}\\
  \GelarPembimbing\\
  \NIDNPembimbing
\end{minipage}
\end{center}

\vspace{2\baselineskip}
\begin{center}
\begin{minipage}[t]{0.6\textwidth}\centering
  Mengetahui,\\[1\baselineskip]
  Ketua Program Studi,\\[3\baselineskip]
  \uline{\NamaKetuaProgramStudi}\\
  \GelarKetuaProgramStudi\\
  \NIDNKetuaProgramStudi
\end{minipage}
\end{center}
\end{singlespace}

%% ===========================================================================
%%  frontmatter/kata-pengantar.tex — Kata Pengantar
%% ===========================================================================
\babfrontmatter{Kata Pengantar}

\begin{singlespace}
% TODO: tulis kata pengantar di sini. Aturan isi: alinea 4-5 kalimat, satu
% gagasan per alinea, istilah asing italic dengan \asing{...}.
% Sebutkan: tujuan penyusunan laporan, syukur, dan terima kasih kepada dosen
% pembimbing (\NamaPembimbing), Ketua Program Studi (\NamaKetuaProgramStudi),
% pimpinan \NamaKampus, pembimbing lapangan dan mentor di \InstansiMagang,
% rekan kerja, dan keluarga.
% Tutup dengan harapan laporan bermanfaat serta permintaan kritik dan saran.

\vspace{2\baselineskip}
\begin{flushright}
\Kota, \TanggalKataPengantar\\[2\baselineskip]
Penulis
\end{flushright}
\end{singlespace}

%% ===========================================================================
%%  frontmatter/abstrak.tex — Abstrak
%%  Abstrak 200-300 kata. Di template ini abstrak memakai spasi 1; ganti
%%  \begin{singlespace} menjadi \begin{spacing}{2} bila pembimbing meminta
%%  spasi 2.
%% ===========================================================================
\babfrontmatter{Abstrak}

\begin{singlespace}
Penulisan karya ilmiah dengan bantuan \asing{large language model} (LLM)
menghadapi dua kelemahan yang berulang: sitasi yang dikarang dan penyuntingan
naskah yang berjalan tanpa sepengetahuan penulis. Laporan magang ini
menyajikan perancangan dan implementasi sistem penulisan riset berbantuan AI
yang dijalankan secara lokal pada komputer peneliti. Jawaban sistem hanya
bersumber pada dokumen yang dikumpulkan peneliti, sitasi hanya sah bila
terverifikasi terhadap pustaka lokal, dan draf tidak pernah masuk ke naskah
tanpa persetujuan manusia.

Sistem memadukan antarmuka web React, layanan \asing{backend} FastAPI dengan
sebelas perute API, basis data SQLite untuk data terstruktur, dan penyimpanan
vektor Chroma untuk potongan sumber setiap proyek. Dokumen sumber dicacah,
diubah menjadi vektor \asing{embedding}, dan disimpan, lalu dipakai pada tahap
pengambilan konteks. Percakapan berjalan dalam tiga mode: \asing{Ask} menjawab
pertanyaan
dari potongan sumber yang ditemukan dan menandai jawaban yang tidak berdasar,
\asing{Plan} hanya menyusun kerangka tanpa menulis naskah, dan \asing{Build}
menyusun draf bagian dari sumber. Konteks yang dikirim ke model dibatasi pada
potongan hasil pengambilan, ringkasan memori percakapan, dan konteks yang
disematkan peneliti. Alur revisi agen memetakan instruksi ke bagian naskah
melalui judul, alias, dan batas kata, lalu menggagalkan instruksi yang ambigu
alih-alih menebak. Setiap sitasi pada teks hasil revisi wajib terverifikasi
pada pustaka lokal yang bersumber dari OpenAlex, sehingga sitasi yang gagal
ditolak atau ditandai. Peneliti memutuskan menerima atau membuang setiap
usulan revisi. Naskah akhir diekspor sebagai berkas LaTeX beserta daftar
pustakanya dan disinkronkan dua arah dengan repositori GitHub untuk setiap
proyek. Pengujian dirancang dalam tiga skenario: keberpijakan sitasi,
ketepatan penargetan bagian naskah, dan ketepatan putaran ekspor serta
sinkronisasi.
% TODO: tambahkan angka hasil pengujian bila sudah tersedia.

\vspace{1\baselineskip}
\textbf{Kata kunci:} penulisan riset berbantuan AI, \asing{retrieval-augmented
generation}, verifikasi sitasi, \asing{human-in-the-loop}, OpenAlex,
\asing{large language model}
\end{singlespace}

%% frontmatter/daftar-isi.tex — Daftar Isi (justify, spasi 1,5)
\daftarisi


%% frontmatter/daftar-tabel.tex — Daftar Tabel (justify, spasi 1,5)
%% Label "Tabel 4.1" berasal dari float tabel yang ditulis di bab.
\daftartabel


%% frontmatter/daftar-grafik.tex — Daftar Grafik (justify, spasi 1,5)
%% Grafik ditulis dengan float khusus agar punya daftar sendiri:
%%   \begin{grafik}[htbp]
%%     \centering
%%     \includegraphics[width=0.8\textwidth]{assets/nama-file.png}
%%     \caption{Judul Grafik}
%%     \label{grafik:contoh}
%%   \end{grafik}
\daftargrafik


%% frontmatter/daftar-gambar.tex — Daftar Gambar (justify, spasi 1,5)
%% Gambar (screenshot, foto, diagram statis) memakai environment figure.
\daftargambar


%% ===========================================================================
%%  frontmatter/daftar-singkatan.tex — Daftar Singkatan dan Akronim
%%  Spasi 1,5, urut abjad. Format entri: \singkatan{SINGKATAN}{Keterangan}
%% ===========================================================================
\babfrontmatter{Daftar Singkatan dan Akronim}

\begin{daftarsingkatan}
% TODO: isi daftar singkatan dan akronim, urut abjad. Contoh format:
% \singkatan{NIM}{Nomor Induk Mahasiswa}
% \singkatan{STMIK}{Sekolah Tinggi Manajemen Informatika dan Komputer}
\end{daftarsingkatan}


%% frontmatter/daftar-lampiran.tex — Daftar Lampiran (justify, spasi 1,5)
%% Terisi otomatis dari \lampiranentry{...} di backmatter/lampiran.tex.
\daftarlampiran



%% ---------------------------------------------------- BAGIAN UTAMA (Arab)
\pagenumbering{arabic}
%% ===========================================================================
%%  Bab 1 Pendahuluan
%%  Aturan isi: alinea 4-5 kalimat, satu gagasan per alinea, spasi 2, awal
%%  alinea menjorok otomatis 1,5 cm, rata kiri-kanan. Istilah asing italic:
%%  \asing{open source}.
%% ===========================================================================
\chapter{Pendahuluan}

\section{Latar Belakang}

Kecerdasan buatan berbasis \asing{large language model} telah mengubah cara
peneliti mencari dan menyusun bahan tulisan. Model tersebut mampu merangkum
puluhan halaman dalam hitungan detik dan menawarkan susunan argumen yang belum
terpikirkan oleh penulis. Kemampuan itu membuat model bahasa besar banyak
dipakai sebagai pendamping dalam penyusunan karya ilmiah, mulai dari pencarian
rujukan sampai penyuntingan kalimat. Kemudahan tersebut, bagaimanapun, membawa
persoalan baru yang tidak selalu tampak pada hasil akhir.

Persoalan pertama adalah sitasi yang tidak dapat dipertanggungjawabkan. Model
bahasa besar menyusun jawaban berdasarkan pola yang dipelajarinya, sehingga
rujukan yang dihasilkan dapat tampak meyakinkan tanpa benar-benar merujuk pada
karya yang ada. \citet{ji2023survey} mengelompokkan kekeliruan semacam itu
sebagai halusinasi, yaitu keluaran yang lancar dibaca tetapi tidak sesuai dengan
sumber maupun fakta. \citet{walters2023fabrication} menguji sitasi yang
dihasilkan salah satu model percakapan dan menemukan bahwa sebagian besar di
antaranya tidak akurat, baik karena karyanya tidak ditemukan maupun karena
rincian penulis, tahun, dan terbitannya keliru. Temuan tersebut menunjukkan
bahwa kelancaran bahasa tidak menjamin ketepatan rujukan.

Persoalan kedua menyangkut kendali penulis atas naskahnya sendiri. Alat bantu
menulis yang bekerja langsung pada dokumen dapat mengganti kalimat tanpa
memberi tahu bagian mana yang berubah. Penulis kemudian harus membandingkan
naskah lama dan naskah baru secara manual untuk mengetahui letak perubahannya.
\citet{amershi2019guidelines} menekankan bahwa sistem berbantuan kecerdasan
buatan perlu bersikap terbuka terhadap tindakan pengguna dan menyediakan cara
yang mudah untuk membatalkan perubahan. Tanpa prinsip tersebut, bantuan yang
dimaksudkan mempercepat pekerjaan justru menambah beban pemeriksaan.

Pendekatan yang banyak dipakai untuk menekan halusinasi adalah
\asing{retrieval-augmented generation}. \citet{lewis2020retrieval} merancang
pendekatan ini dengan mengambil dokumen yang relevan terlebih dahulu, lalu
menyusun jawaban berdasarkan dokumen tersebut. \citet{gao2023retrieval}
menjelaskan bahwa mutu jawaban pada pendekatan ini sangat ditentukan oleh
ketepatan tahap pengambilan sumber. Namun, sebagian besar implementasi yang
tersedia masih menempatkan dokumen peneliti pada layanan awan dan belum
membedakan tahap perencanaan dari tahap penulisan naskah. Pemisahan tahap itu
penting agar peneliti dapat menilai rencana tulisan sebelum ada kalimat yang
masuk ke dalam dokumen.

Kebutuhan tersebut muncul secara nyata di Broadband Multimedia Wireless
Laboratory, National Taiwan University of Science and Technology. Laboratorium
ini mengerjakan penelitian di bidang jaringan nirkabel dan multimedia, sehingga
setiap anggota tim mengelola banyak artikel, laporan, dan catatan percobaan
secara bersamaan. Peneliti perlu menelusuri kembali sumber yang pernah dibaca,
memastikan bahwa rujukan yang dipakai benar-benar ada, dan menjaga agar naskah
yang sedang disusun tidak berubah tanpa sepengetahuan mereka. Pekerjaan
penelusuran semacam itu menyita waktu yang seharusnya dapat dipakai untuk
menguji gagasan dan menganalisis hasil.

Berdasarkan keadaan tersebut, selama pelaksanaan magang dirancang dan
diimplementasikan sebuah sistem pencarian data riset berbantuan kecerdasan
buatan yang dijalankan pada komputer peneliti. Sistem ini mencari bagian
dokumen yang relevan dengan pertanyaan peneliti, lalu menyusun jawaban hanya
dari bagian yang ditemukan. Sistem menyediakan tiga cara kerja, yaitu menjawab
pertanyaan, menyusun kerangka, dan menyusun draf bagian naskah. Setiap usulan
perubahan harus disetujui peneliti sebelum dimasukkan ke dalam naskah, dan
setiap sitasi diperiksa terhadap pustaka yang tersimpan pada proyek.

Sistem yang dibangun juga mencatat asal setiap potongan sumber sehingga jawaban
dapat ditelusuri kembali. Pencatatan ini penting karena peneliti tidak cukup
mempercayai jawaban yang terbaca lancar; mereka perlu mengetahui dari dokumen
mana sebuah pernyataan berasal. Naskah yang telah disetujui diekspor dalam
format LaTeX beserta daftar pustakanya, lalu disinkronkan dengan repositori
proyek pada GitHub. Dengan demikian, riwayat perubahan naskah tetap terekam dan
dapat diperiksa oleh pembimbing maupun rekan tim.

Laporan magang ini merekam perancangan, implementasi, dan evaluasi sistem
tersebut sebagai bagian dari pelaksanaan magang pada periode 12 Januari sampai
28 September 2026. Laporan disusun mengikuti Pedoman Penulisan Laporan Magang
STMIK Kuwera agar dapat dipertanggungjawabkan secara akademik. Susunan bab
disesuaikan dengan urutan yang ditetapkan pedoman, mulai dari dasar pemikiran
sampai kesimpulan. Selain menjadi dokumen pertanggungjawaban, laporan ini
diharapkan menjadi rujukan bagi mahasiswa yang akan mengerjakan pekerjaan
serupa pada periode berikutnya.

\section{Tujuan Magang}

Magang dilaksanakan untuk menghubungkan pengetahuan yang diperoleh di bangku
kuliah dengan pekerjaan nyata di lingkungan laboratorium riset. Penulis
ditempatkan pada pekerjaan pengembangan perangkat lunak yang menopang kegiatan
penelitian, sehingga dapat mengamati langsung bagaimana kebutuhan peneliti
diterjemahkan menjadi fitur sistem. Penempatan tersebut juga memberi kesempatan
kepada penulis untuk menguji keputusan rancangan pada pengguna yang sebenarnya.
Tujuan umum tersebut dijabarkan menjadi beberapa tujuan khusus sebagai berikut.

\begin{enumerate}
  \item Merancang sistem pencarian data riset berbantuan kecerdasan buatan yang
        dijalankan secara lokal pada komputer peneliti.
  \item Mengimplementasikan pencarian dokumen sumber sehingga jawaban yang
        dihasilkan dapat ditelusuri kembali kepada bagian asalnya.
  \item Mengimplementasikan verifikasi sitasi terhadap pustaka lokal yang
        tersimpan pada setiap proyek penelitian.
  \item Mengimplementasikan alur revisi yang menempatkan persetujuan peneliti
        sebagai syarat sebelum naskah berubah.
  \item Menguji ketepatan sumber, ketepatan penentuan bagian naskah, serta
        keberhasilan ekspor dan sinkronisasi naskah.
\end{enumerate}

Melalui kelima tujuan tersebut, magang diarahkan pada satu hasil yang utuh,
yaitu perangkat lunak yang dapat dipakai dalam pekerjaan penulisan sehari-hari.
Setiap tujuan memiliki ukuran keberhasilan yang diperiksa pada tahap pengujian.
Penulis juga mencatat kendala yang muncul selama pengerjaan agar dapat dijadikan
bahan perbaikan pada pengembangan berikutnya.

\section{Manfaat Magang}

Pelaksanaan magang memberikan manfaat bagi tiga pihak yang terlibat, yaitu
peserta magang, program studi, dan instansi tujuan magang. Manfaat tersebut
tidak berdiri sendiri, melainkan saling berkaitan sebagai hasil dari pekerjaan
yang sama. Peserta magang memperoleh pengalaman kerja, program studi memperoleh
contoh penerapan ilmu, dan instansi memperoleh perangkat yang membantu
pekerjaannya. Uraian berikut menjelaskan bentuk manfaat pada masing-masing
pihak.

\subsection{Manfaat bagi Peserta Magang}

Manfaat utama bagi peserta magang adalah pengalaman mengerjakan perangkat
lunak dari tahap perancangan sampai pengujian. Penulis belajar menerjemahkan
kebutuhan pengguna yang masih umum menjadi rancangan fitur yang dapat
dikerjakan. Penulis juga berlatih menulis kode, membaca kesalahan, dan
memperbaiki sistem yang sudah berjalan tanpa mengganggu data pengguna.
Pengalaman tersebut tidak diperoleh hanya dari membaca buku, karena setiap
keputusan harus diuji pada pemakaian yang sebenarnya.

Selain keterampilan teknis, magang melatih cara berkomunikasi dengan peneliti
yang menjadi pengguna sistem. Penulis belajar menanyakan kebutuhan dengan
pertanyaan yang tepat dan menjelaskan keputusan teknis tanpa memakai istilah
yang membingungkan. Pengalaman ini menumbuhkan kebiasaan mendokumentasikan
pekerjaan agar dapat dilanjutkan oleh orang lain. Penulis juga memperoleh
gambaran mengenai budaya kerja penelitian di laboratorium universitas luar
negeri.

\subsection{Manfaat bagi Program Studi}

Bagi Program Studi Sistem Informasi, magang ini menghasilkan contoh penerapan
ilmu sistem informasi pada pekerjaan penelitian. Kasus yang dikerjakan
memperlihatkan bahwa analisis kebutuhan, perancangan basis data, dan
pengembangan antarmuka tidak hanya dipakai pada sistem bisnis, tetapi juga pada
perangkat pendukung riset. Contoh tersebut dapat dipakai sebagai bahan
pembelajaran pada mata kuliah yang relevan. Dengan begitu, materi kuliah
memperoleh gambaran penerapan yang dekat dengan pekerjaan lulusan.

Magang ini juga memperluas jejaring kerja sama program studi dengan institusi
pendidikan di luar negeri. Hubungan yang terbentuk selama periode magang dapat
dimanfaatkan untuk kegiatan pertukaran mahasiswa maupun kerja sama penelitian
pada masa berikutnya. Pengalaman penulis selama mengikuti program tersebut
dapat dibagikan kepada mahasiswa lain sebagai gambaran dunia kerja di bidang
tersebut. Program studi memperoleh masukan mengenai kemampuan yang perlu
diperkuat pada mahasiswa yang berminat bekerja di bidang riset.

\subsection{Manfaat bagi Instansi Tujuan Magang}

Bagi Broadband Multimedia Wireless Laboratory, magang ini menghasilkan
perangkat lunak yang membantu anggota laboratorium mengelola bahan penelitian.
Sistem yang dibangun menyimpan dokumen sumber pada komputer peneliti dan
memungkinkan penelusuran kembali bagian yang relevan dengan pertanyaan. Dengan
demikian, waktu yang biasa dipakai untuk mencari ulang dokumen dapat dialihkan
ke pekerjaan analisis. Perangkat tersebut juga dapat dipakai oleh anggota baru
yang sedang mempelajari topik penelitian laboratorium.

Laboratorium juga memperoleh alur kerja yang mencatat persetujuan peneliti
sebelum naskah berubah. Pencatatan tersebut memudahkan peninjauan bersama antara
anggota tim dan pembimbing, karena setiap perubahan dapat ditelusuri. Selain
itu, sistem menyimpan naskah beserta daftar pustakanya pada repositori proyek,
sehingga riwayat penyusunan dokumen tetap tersedia setelah mahasiswa magang
menyelesaikan periodenya. Hasil kerja ini menjadi titik awal bagi pengembangan
perangkat serupa pada proyek penelitian berikutnya.


%% ===========================================================================
%%  Bab 2 Dasar Pemikiran
%% ===========================================================================
\chapter{Dasar Pemikiran}

\section{Dasar Pemikiran Pemilihan Lokasi Magang}
% TODO: alasan memilih \InstansiMagang{} sebagai lokasi magang, misalnya
% kesesuaian dengan bidang ilmu, ketersediaan laboratorium, dan peluang kerja
% sama internasional.

\section{Dasar Pemikiran Pemilihan Bidang Magang}
% TODO: alasan memilih bidang penulisan riset berbantuan AI, misalnya
% relevansi dengan \ProgramStudi, kebutuhan instansi, dan bekal karier.

%% ===========================================================================
%%  Bab 3 Profil Organisasi dan Aktivitas Magang
%% ===========================================================================
\chapter{Profil Organisasi dan Aktivitas Magang}

\section{Gambaran Umum Organisasi}
% TODO: profil \InstansiMagang: sejarah singkat, bidang kerja, struktur
% organisasi, dan penempatan unit tempat magang.

\section{Aktivitas Magang}
% TODO: uraikan aktivitas selama magang, boleh dibagi per pekan, per proyek,
% atau per tahap pekerjaan. Sebutkan periode \PeriodeMagang.

\section{Metode Pelaksanaan Magang}
% TODO: cara pelaksanaan kerja, misalnya daring/luring, bimbingan berkala,
% rapat tim, dan alur pelaporan.

\section{Bentuk-bentuk Dukungan}
% TODO: satu alinea pengantar, lalu rincian pada tiga subbagian di bawah.

\subsection{Dukungan dari Mahasiswa}
% TODO: peran, inisiatif, dan kesiapan mahasiswa.

\subsection{Dukungan dari Institusi Tujuan Magang}
% TODO: fasilitas, mentor, data, dan akses yang diberikan instansi.

\subsection{Dukungan dari Program Studi}
% TODO: dukungan administratif, pembimbing, dan kebijakan kampus.

\section{Kendala-kendala}
% TODO: satu alinea pengantar, lalu rincian pada tiga subbagian di bawah.

\subsection{Kendala dari Mahasiswa}
% TODO: kendala teknis dan nonteknis dari sisi mahasiswa.

\subsection{Kendala dari Institusi Tujuan Magang}
% TODO: kendala dari sisi instansi, misalnya akses data atau perizinan.

\subsection{Kendala dari Program Studi}
% TODO: kendala dari sisi kampus, misalnya administrasi atau jadwal.

%% ===========================================================================
%%  Bab 4 Pembahasan
%%  Bab terpanjang: isi teori, analisis, dan rekomendasi perbaikan.
%%  Contoh penulisan tabel, gambar, grafik, dan kode ada di bawah (komentar).
%% ===========================================================================
\chapter{Pembahasan}

\section{Landasan Teori}
% TODO: teori yang menopang pekerjaan magang. Sitasi memakai \citep{kunci}
% atau \citet{kunci} dengan entri di backmatter/references.bib.

\section{Analisis}
% TODO: analisis kebutuhan, kondisi sistem sebelum/sesudah, dan hasil kerja.

\section{Rekomendasi Perbaikan}
% TODO: usulan perbaikan untuk instansi, sistem, atau proses kerja berikutnya.

%% --------------------------------------------------------------------------
%%  Contoh tabel (judul tabel di atas tabel, nomor per bab: Tabel 4.1)
%% --------------------------------------------------------------------------
%% \begin{table}[htbp]
%%   \centering
%%   \caption{Judul Tabel yang Ditulis Capitalize Each Word}
%%   \label{tabel:contoh}
%%   \begin{tabular}{@{}lcc@{}}
%%     \toprule
%%     Kolom A & Kolom B & Kolom C \\
%%     \midrule
%%     Isi & Isi & Isi \\
%%     \bottomrule
%%   \end{tabular}
%% \end{table}

%% --------------------------------------------------------------------------
%%  Contoh gambar (judul gambar di bawah gambar: Gambar 4.1)
%% --------------------------------------------------------------------------
%% \begin{figure}[htbp]
%%   \centering
%%   \includegraphics[width=0.8\textwidth]{assets/nama-file.png}
%%   \caption{Judul Gambar}
%%   \label{gambar:contoh}
%% \end{figure}

%% --------------------------------------------------------------------------
%%  Contoh grafik (masuk ke Daftar Grafik: Grafik 4.1)
%% --------------------------------------------------------------------------
%% \begin{grafik}[htbp]
%%   \centering
%%   \includegraphics[width=0.8\textwidth]{assets/nama-file.png}
%%   \caption{Judul Grafik}
%%   \label{grafik:contoh}
%% \end{grafik}

%% --------------------------------------------------------------------------
%%  Contoh kutipan kode program (di dalam bab)
%% --------------------------------------------------------------------------
%% \begin{singlespace}
%% \begin{lstlisting}[language=Python,caption={Contoh Kode Program}]
%% def contoh():
%%     return "halo"
%% \end{lstlisting}
%% \end{singlespace}

%% ===========================================================================
%%  Bab 5 Refleksi Diri
%%  Tiga subbab di bawah mengikuti butir refleksi pada pedoman.
%% ===========================================================================
\chapter{Refleksi Diri}

\section{Manfaat Magang untuk Hardskill, Softskill, Kemampuan Kognitif, dan Aktualisasi Diri}
% TODO: tuliskan kemampuan teknis, kemampuan bekerja sama, cara berpikir, dan
% pengalaman mengaktualisasikan diri selama magang.

\section{Kunci Sukses Bekerja Berdasarkan Pengalaman Magang}
% TODO: kebiasaan atau sikap yang membuat pekerjaan berhasil, misalnya
% komunikasi, disiplin, dokumentasi, dan manajemen waktu.

\section{Rencana Perbaikan atau Pengembangan Diri}
% TODO: rencana belajar dan perbaikan diri setelah magang.

%% ===========================================================================
%%  Bab 6 Kesimpulan
%% ===========================================================================
\chapter{Kesimpulan}

\section{Kesimpulan}
% TODO: jawaban atas rumusan masalah dan tujuan magang di Bab 1.

\section{Saran}
% TODO: saran untuk instansi, kampus, dan mahasiswa yang akan magang.



%% ----------------------------------- BAGIAN AKHIR (melanjutkan angka Arab)
%% ===========================================================================
%%  backmatter/daftar-pustaka.tex — Daftar Pustaka
%%  Gaya sitasi nama-tahun (apalike). Sumber ditulis di references.bib.
%% ===========================================================================
\clearpage

\begingroup
\setstretch{1.5}
\renewcommand{\bibname}{Daftar Pustaka}
\addcontentsline{toc}{chapter}{Daftar Pustaka}
\bibliographystyle{apalike}
\bibliography{backmatter/references}
\endgroup


%% ===========================================================================
%%  backmatter/lampiran.tex — Lampiran
%%  Nomor halaman melanjutkan angka Arab. Setiap lampiran diberi judul
%%  "Lampiran <n>. <Judul>" lewat \lampiranentry{...} agar otomatis masuk
%%  ke Daftar Lampiran.
%% ===========================================================================
\babfrontmatter{Lampiran}

% TODO: buat lampiran sesuai kebutuhan. Format:
%
%   \lampiranentry{Lampiran 1. Kode Program Utama}
%   \begin{singlespace}
%   \begin{lstlisting}[language=Python]
%   ...
%   \end{lstlisting}
%   \end{singlespace}
%
%   \lampiranentry{Lampiran 2. Surat Keterangan Magang}
%   \includegraphics[width=\textwidth]{assets/surat-magang.png}



\end{document}

